% 第 1 章：集合与常用逻辑用语
\chapter{集合与常用逻辑用语}

\section{基础题}

\begin{problem}\level{易}
已知 $A=\{1,2,3\}$，$B=\{2,3,4\}$，则 $A\cap B=$？
\end{problem}
\begin{solution}
$A\cap B=\{2,3\}$。
\end{solution}

\begin{problem}\level{易}
设集合 $A=\{x\mid x^2-3x+2=0\}$，则 $A$ 的子集个数为？
\end{problem}
\begin{solution}
解方程得 $x=1$ 或 $x=2$，故 $A=\{1,2\}$，子集个数为 $2^2=4$。
\end{solution}

\begin{problem}\level{中}
已知集合 $A=\{x\mid -1<x\le 2\}$，$B=\{x\mid x\le 0\}$，求 $A\cup B$ 与 $A\cap \complement_{\mathbb{R}}B$。
\end{problem}
\begin{solution}
$A\cup B=\{x\mid x\le 2\}$；$\complement_{\mathbb{R}}B=\{x\mid x>0\}$，故 $A\cap\complement_{\mathbb{R}}B=\{x\mid 0<x\le 2\}$。
\end{solution}

\begin{problem}\level{中}
设 $a,b\in\mathbb{R}$，集合 $\{1,a+b,a\}=\left\{0,\dfrac{b}{a},b\right\}$，求 $b-a$ 的值。
\end{problem}
\begin{solution}
由互异性知 $a\ne 0$，故 $a+b=0$，即 $b=-a$。于是 $\dfrac{b}{a}=-1$，比较得 $b=1$，$a=-1$，所以 $b-a=2$。
\end{solution}

\begin{problem}\level{难}
设集合 $M=\left\{x\mid x=\dfrac{k}{2}+\dfrac{1}{4},k\in\mathbb{Z}\right\}$，$N=\left\{x\mid x=\dfrac{k}{4}+\dfrac{1}{2},k\in\mathbb{Z}\right\}$，试判断 $M$ 与 $N$ 的关系。
\end{problem}
\begin{solution}
$M=\left\{\dfrac{2k+1}{4}\mid k\in\mathbb{Z}\right\}$ 为所有奇数除以 $4$；$N=\left\{\dfrac{k+2}{4}\mid k\in\mathbb{Z}\right\}$ 为所有整数除以 $4$。因此 $M\subsetneqq N$。
\end{solution}

\section{量词与解集结构}

\begin{problem}\level{易}
写出“$\exists x\in\mathbb{R}$，$x^2+1<0$”的否定，并判断原命题的真假。
\end{problem}
\begin{solution}
否定为“$\forall x\in\mathbb{R}$，$x^2+1\geq0$”。原命题为假。
\end{solution}

\begin{problem}\level{易}
已知 $f(x)=x^2+1$，把“$f(x)>0$ 对任意 $x\in\mathbb{R}$ 成立”写成符号形式，并用集合语言说明。
\end{problem}
\begin{solution}
$\forall x\in\mathbb{R}$，$f(x)>0$。集合语言：$\{x\mid f(x)>0\}=\mathbb{R}$（解集等于定义域）。
\end{solution}

\begin{problem}\level{中}
设 $p:\exists x\in\mathbb{R}$，$x^2+ax+1<0$。若 $p$ 为真，求 $a$ 的取值范围。
\end{problem}
\begin{solution}
解集 $\{x\mid x^2+ax+1<0\}$ 非空 $\iff\Delta=a^2-4>0$，即 $a>2$ 或 $a<-2$。
\end{solution}

\begin{problem}\level{中}
设 $A=\{x\mid x^2-2x-3<0\}$，$B=\{x\mid x>a\}$。若 $\forall x\in A$ 都有 $x\in B$，求 $a$ 的取值范围。
\end{problem}
\begin{solution}
$A=(-1,3)$，$B=(a,+\infty)$。$\forall x\in A,\ x\in B$ 即 $A\subseteq B$，故 $a\le-1$。
\end{solution}

\begin{problem}\level{中}
用反证法证明：对任意 $x\in\mathbb{R}$，都有 $x^2+x+1>0$。
\end{problem}
\begin{solution}
假设存在 $x$ 使 $x^2+x+1\leq0$，则 $\left(x+\dfrac12\right)^2+\dfrac34\leq0$，
即 $\left(x+\dfrac12\right)^2\leq-\dfrac34$，与平方非负矛盾。
\end{solution}

\begin{problem}\level{中}
设集合 $S$ 满足：$1\in S$，且若 $a\in S$，则 $\dfrac{1}{1-a}\in S$。已知 $2\in S$，写出 $S$ 的全部元素。
\end{problem}
\begin{solution}
由 $2\in S$ 得 $\dfrac{1}{1-2}=-1\in S$；由 $-1\in S$ 得 $\dfrac{1}{1-(-1)}=\dfrac12\in S$；
由 $\dfrac12\in S$ 得 $\dfrac{1}{1-\frac12}=2\in S$（回到起点）。
故 $S=\left\{2,-1,\dfrac12\right\}$。
\end{solution}

\begin{problem}\level{难}
证明：不存在长度大于 $0$ 的区间 $I$，使 $\sin x=0$ 对所有 $x\in I$ 成立。
\end{problem}
\begin{solution}
$\sin x=0$ 的解集为 $\{k\pi\mid k\in\mathbb{Z}\}$，是离散点集，长度为 $0$。
任何长度大于 $0$ 的区间都不可能含于其中，故不存在。
\end{solution}

\section{章节综合练习}

\begin{problem}\level{易}
用描述法表示集合 $\{1,3,5,7,9\}$。
\end{problem}
\begin{solution}
$\{x\mid x=2n-1,\ n\in\mathbb{N}^*,\ n\le 5\}$。
\end{solution}

\begin{problem}\level{中}
已知 $A=\{x\mid x^2-5x+6<0\}$，$B=\{x\mid 1<x<a\}$，若 $A\subseteq B$，求 $a$ 的取值范围。
\end{problem}
\begin{solution}
$A=(2,3)$。由 $A\subseteq B$ 得 $a\ge 3$。
\end{solution}

\begin{problem}\level{难}
设集合 $A=\{1,2,3,\dots,10\}$，求 $A$ 的所有子集中，元素和为偶数的子集个数。
\end{problem}
\begin{solution}
$A$ 中有 $5$ 个奇数、$5$ 个偶数。子集元素和的奇偶性由所含奇数个数决定，故需取偶数个奇数：$\binom{5}{0}+\binom{5}{2}+\binom{5}{4}=16$ 种；偶数可任取，共 $2^5=32$ 种。合计 $16\times 32=512$。
\end{solution}
